% !TeX program = lualatex
% !TeX encoding = utf8
% !TeX spellcheck = uk_UA

\documentclass{article}
\usepackage[fontsize=14pt]{fontsize}

\usepackage{fontsetup}
\usepackage[english, russian, ukrainian]{babel}
\usepackage{microtype}


\usepackage[%
	a4paper,%
	footskip=1cm,%
	headsep=0.3cm,%
	top=2cm, %поле сверху
	bottom=2cm, %поле снизу
	left=2cm, %поле ліворуч
	right=2cm, %поле праворуч
    ]{geometry}

\usepackage{amsmath}
\usepackage{graphicx}
\usepackage{floatflt}

\usepackage[%colorlinks=true,
	%urlcolor = blue, %Colour for external hyperlinks
	%linkcolor  = malina, %Colour of internal links
	%citecolor  = green, %Colour of citations
	bookmarks = true,
	bookmarksnumbered=true,
	unicode,
	linktoc = all,
	hypertexnames=false,
	pdftoolbar=false,
	pdfpagelayout=TwoPageRight,
	pdfauthor={Ponomarenko S.M. aka sergiokapone},
	pdfdisplaydoctitle=true,
	pdfencoding=auto
	]%
	{hyperref}



\title{}
\date{}

\begin{document}


\textbf{Модель.}
Для плоского або не\-плоского Всесвіту з компонентами
$(\Omega_m, \Omega_\Lambda, \Omega_k = 1 - \Omega_m - \Omega_\Lambda - \Omega_r)$
безрозмірна функція Хаббла:
\begin{equation}
  E(z) = \sqrt{\Omega_r(1+z)^4 + \Omega_m(1+z)^3
               + \Omega_k(1+z)^2 + \Omega_\Lambda}.
\end{equation}
Супутня відстань:
\begin{equation}
  d_C(z) = \frac{c}{H_0}\int_0^z \frac{dz'}{E(z')}.
\end{equation}
Усереднена BAO-відстань:
\begin{equation}
  D_V(z) \equiv \left[z\,d_C^2(z)\,\frac{c}{H(z)}\right]^{1/3},
\end{equation}
де $H(z) = H_0 E(z)$.
Спостережна величина~--- відношення $D_V(z)/r_s$,
де $r_s = 147.09$~Мпк~--- звуковий горизонт (Planck~2018).

\textbf{Підбір параметрів.}
Мінімізується:
\begin{equation}
  \chi^2(H_0, \Omega_m, \Omega_\Lambda)
  = \sum_{i=1}^{N}
    \left(\frac{D_V^{\rm obs}(z_i)/r_s
               - D_V^{\rm th}(z_i)/r_s}{\sigma_i}\right)^2.
\end{equation}
Три вільні параметри: $H_0$, $\Omega_m$, $\Omega_\Lambda$.
Параметр кривини не фіксується: $\Omega_k = 1 - \Omega_m - \Omega_\Lambda - \Omega_r$.


\end{document}


